% BHU PYQ -- Professional Exam Template
\documentclass[11pt,a4paper]{article}

\usepackage[a4paper,top=2.5cm,bottom=2.5cm,left=2.8cm,right=2.8cm,headheight=14pt]{geometry}
\usepackage{amsmath,amssymb}
\usepackage[T1]{fontenc}
\usepackage[utf8]{inputenc}
\usepackage{lmodern}
\usepackage{microtype}
\usepackage{fancyhdr}
\usepackage{enumitem}
\usepackage{listings}
\usepackage{hyperref}
\usepackage{lastpage}
\usepackage{parskip}

\hypersetup{colorlinks=true,urlcolor=black,linkcolor=black,
  pdftitle={BCS-503: Discrete Mathematical Structures -- BHU PYQ}}

\pagestyle{fancy}
\fancyhf{}
\renewcommand{\headrulewidth}{0.4pt}
\renewcommand{\footrulewidth}{0.4pt}
\fancyhead[L]{\small   \textit{Banaras Hindu University}}
\fancyhead[R]{\small   \textit{BCS-503: Discrete Math Structures}}
\fancyfoot[C]{\small Page  \thepage\ of \pageref{LastPage}}

\lstset{
  basicstyle=\small tfamily,
  frame=single,
  breaklines=true,
  columns=flexible,
  keepspaces=true
}


\newlist{parts}{enumerate}{2}
\setlist[parts,1]{label=(\alph*),leftmargin=*,itemsep=4pt,topsep=3pt}
\setlist[parts,2]{label=(\roman*),leftmargin=*,itemsep=2pt,topsep=2pt}

\newcommand{\pts}[1]{\hfill   \textit{[#1]}}


\begin{document}

\begin{center}
  {\large\bfseries BANARAS HINDU UNIVERSITY}\\[2pt]
  {
\normalsize B.Sc. (Hons.) Semester V Examination, 2016-17}\\[6pt]
  \rule{0.8\linewidth}{0.5pt}\\[6pt]
  {\large\bfseries Paper: BCS-503 --- Discrete Mathematical Structures}\\[4pt]
  {
\normalsize Time: Three Hours \hfill Full Marks: 70}
\end{center}

\rule{\linewidth}{0.4pt}

\smallskip



\noindent   \textit{Instructions: Answer any five questions. Terms/Symbols used have their standard meanings.}

\bigskip

\medskip


\noindent   \textbf{Question 1.} Define the following terminologies and explain with appropriate examples: \pts{$2 imes7=14$}

\begin{parts}
  \item Poset
  \item Regular Graph
  \item Clique
  \item Fundamental Circuit
  \item Bipartite Graph
  \item Minimal Covering
  \item De Morgan's Law
\end{parts}

\medskip


\noindent   \textbf{Question 2.}

\begin{parts}
  \item Explain partition of a set and equivalence classes with a suitable example. \pts{3}
  \item What is a Hasse Diagram? Draw the Hasse Diagram for the relation $R$ on set $A = \{a, b, c, d, e\}$ where $R$ is given as:
  \[ R = \{\langle a, a \rangle, \langle a, c \rangle, \langle a, d \rangle, \langle a, e \rangle, \langle b, b \rangle, \langle b, c \rangle, \langle b, d \rangle, \langle b, e \rangle, \langle c, c \rangle, \langle c, d \rangle, \langle c, e \rangle, \langle d, d \rangle, \langle e, e \rangle\} \] \pts{4}
  \item Find the inverse of the function $f(x) = 4e^{x+2}$. \pts{3}
  \item A mobile manufacturing company assembled $30$ mobiles. The features available are Bluetooth, Wi-Fi and 4G network. It is known that $15$ of the mobiles have Bluetooth, $8$ of them have Wi-Fi, $6$ of them have 4G network, and $3$ of them have all three features. Find at least how many mobiles do not have any features at all. \pts{4}
\end{parts}

\medskip


\noindent   \textbf{Question 3.}

\begin{parts}
  \item For an algebraic system $(A, \lor, \land)$ defined by a lattice $(A, \le)$, show that: \pts{5}
  
\begin{parts}
    \item $a \le a \lor b$, $a \land b \le a$
    \item $a \lor (a \land b) = a$, $a \land (a \lor b) = a$ where $a, b \in A$.
  \end{parts}
  \item Define Boolean lattice and Boolean algebra. Prove that an element in a Boolean lattice has a unique complement. \pts{5}
  \item Given a function $f: \{0, 1\}^3  o \{0, 1\}$ defined as:
  
\begin{align*}
    f(0,0,0) &= 1, & f(0,0,1) &= 0, \\
    f(0,1,0) &= 1, & f(0,1,1) &= 0, \\
    f(1,0,0) &= 0, & f(1,0,1) &= 0, \\
    f(1,1,0) &= 0, & f(1,1,1) &= 1.
  \end{align*}
  Express the function in disjunctive normal form and conjunctive normal form of Boolean expression. \pts{4}
\end{parts}

\medskip


\noindent   \textbf{Question 4.}

\begin{parts}
  \item Determine whether each of the following compound propositions is a tautology or not: \pts{3}
  
\begin{parts}
    \item $(
eg p \land 
eg q)  o (p  o q)$
    \item $(p  o q)  o (p \land 
eg q)$
    \item $[p \land (p  o 
eg q)]  o q$
  \end{parts}
  \item Show that $
eg (p \lor (
eg p \land q)) \equiv 
eg p \land 
eg q$ and $
eg \forall x (P(x)  o Q(x)) \equiv \exists x (P(x) \land 
eg Q(x))$ without truth table. \pts{4}
  \item Solve the following recurrence relations: \pts{7}
  
\begin{parts}
    \item $4 a_r - 20 a_{r-1} + 17 a_{r-2} - 4 a_{r-3} = 0$
    \item $a_r + 5 a_{r-1} + 6 a_{r-2} = 3^r - 2r + 1$
  \end{parts}
\end{parts}

\medskip


\noindent   \textbf{Question 5.}

\begin{parts}
  \item Prove that the number of vertices of odd degree in a graph is always even. \pts{3}
  \item Draw the graphs of the following types: \pts{5}
  
\begin{parts}
    \item Graph is both Hamiltonian and Eulerian.
    \item Graph is Eulerian but not Hamiltonian.
    \item Graph is Hamiltonian but not Eulerian.
  \end{parts}
  \item What is meant by center of a graph? Prove that a tree may have one or two centers. \pts{2}
  \item What will be the minimum and maximum possible height of a binary tree with $n$ vertices? \pts{2}
  \item What are the Kuratowski's two graphs? Also draw these two graphs. \pts{2}
\end{parts}

\medskip


\noindent   \textbf{Question 6.}

\begin{parts}
  \item How to prove that a graph is non-planar? Show with an example. \pts{4}
  \item What is geometric dual of a graph? If $r$ and $\mu$ denote the rank and nullity of a graph $G$, and if $r^*$ and $\mu^*$ denote the same for the dual graph $G^*$, then prove that $r^* = \mu$ and $\mu^* = r$. \pts{4}
  \item Let $G = (V, E)$ be a connected weighted and undirected graph, where $V = \{a, b, c, d, e, f, g\}$ is the set of vertices and
  \[ E = \{\langle (a, b), 3 \rangle, \langle (a, c), 2 \rangle, \langle (b, d), 1 \rangle, \langle (c, d), 5 \rangle, \langle (c, f), 7 \rangle, \langle (d, e), 3 \rangle, \langle (d, f), 6 \rangle, \langle (e, g), 8 \rangle, \langle (f, g), 4 \rangle\} \]
  is the set of edges along with weights. Find the minimal spanning tree of $G$ using Kruskal's and Prim's algorithms with demonstration of steps in both methods. \pts{6}
\end{parts}

\medskip


\noindent   \textbf{Question 7.}

\begin{parts}
  \item What is chromatic partitioning? Explain with an example. \pts{3}
  \item What is chromatic polynomial? Find the chromatic polynomial of a complete graph of $n$ vertices and a tree of $n$ vertices. \pts{5}
  \item Explain the reduced incidence matrix $A_f$, fundamental circuit matrix $B_f$, and fundamental cutset matrix $C_f$ of a connected graph. Use a common graph and common spanning tree to explain these matrices. What is the relationship among them? \pts{6}
\end{parts}

\vfill
\rule{\linewidth}{0.4pt}

\begin{center}\small
\href{https://bhu-exam.vercel.app/}{Click here to visit our website}\\[4pt]\href{https://me-aryan.vercel.app/}{Click here to visit our contributor}\\[4pt]\href{https://quashan-me.vercel.app/}{Click here to visit our contributor}
\end{center}
\end{document}